\section{Introduction}

The sign problem for the Weil form brings together two
requirements: preservation of all its arithmetic contributions
and control of the form on complete function spaces. The global
criterion relates its positivity to the Riemann hypothesis.
A spectral construction seeking to apply that criterion must
therefore specify the product, weights and domains determining
the form and prove how its cross terms are transmitted. This
article develops that relation from an arithmetic of paths
with memory and obtains exact reconstruction, transport and
coercivity results.

The principal functional contribution is a Schur reduction for
each fixed window. A sufficiently fine partition separates
cell averages from an infinite-dimensional space of zero-mean
details whose form is coercive. Exact elimination of these
details preserves their entire coupling with the averages
and produces a finite matrix with the same negative index
as the complete form. The result supplies an exact criterion,
not a replacement of the function space by samples. Its
application on $(-1/18,1/18)$ proves the lower bound
$\frac12\|f\|_2^2$ on the entire joint readout domain.
A reversible transport extends this bound to complete
collections of test functions with exponential tails.

\subsection*{Arithmetic construction and preserved structure}

The starting point is a formal mathematical core called
Triadic Modular Holography, henceforth HMT. Its object
\emph{All Possible Paths} (APP) combines paths and arithmetic
evaluations on the support $X_9=(\mathbb Z/9\mathbb Z)^2$.
The cardinal generators $\{(1,0),(-1,0),(0,1),(0,-1)\}$
define the Cayley graph of its additive group, a
two-dimensional cyclic lattice with toroidal realization.
Sum and product evaluations act on the same positions,
each with its residue and quotient. The support specifies
where a route moves; the arithmetic sheets specify what
is computed and preserved during that movement.

TRIT is the oriented ternary signature $\{-1,0,+1\}$,
obtained through balanced representatives modulo $3$.
It distinguishes the hyperbolic, parabolic and elliptic
types in its quadratic realization. The
\emph{Topological Phase Kernel}, henceforth TPK, composes
phase selection, transport and memory update. Its state
preserves positions, sheets, orientation, quotients, prefixes
and boundary. Nonadic prolongation transports this
information between depths: phase return advances history
instead of restarting it.

The joint discrete structure of the continuum is developed
on these same states and their compatibilities. Its
correlative constructions preserve restrictions, incidences,
forms and memory within a single refinement system.
The Archimedean readout obtains values from compatible
cylinders, and the boundary readout retains the incidences
used by subsequent realizations. The present arithmetic
and spectral problem preserves this common origin, with
the operations and domains actually entering each result.

The center--boundary relation provides a first instance of
reconstruction. In the multiaffine APP chart, the lifted
boundary determines the interior; residual values are
completed by the quotients distinguishing their integer
values. Continuation signatures require the sheets to be
preserved, and the vacancy calendar retains the discrepancy
between capacity and depth. Section~\ref{sec:centro-frontera}
places the center $(5,5)$, the regional signature $555555$,
the readouts $72,27,77,22$ and the intervals $21/22$ and
$6/7$ within this construction. Their role specifies
transportable information before any spectral identification.

\subsection*{From the product to the completed function}

The native product composes the two APP sheets through a
transducer of digits and quotients. Its evaluation invariant
proves exactness of the operation on nonadic words, and
uniqueness of normal form subsequently identifies integer
multiplication. Product fibres determine its decompositions;
the absence of a nontrivial decomposition selects the
irreducibles. Theorem~\ref{rz-num-01-6} identifies their
evaluations with positive primes.

Factorization and arithmetic clocks describe different
properties of these objects. The former determines their
multiplicative components; the latter records returns of
residual orbits. The mode-occupation representation preserves
the multiplicity of every irreducible and leads to the
Euler product. The trace over all occupations and the
trace over individual modes are related by the prime--power
moments of Proposition~\ref{rz-num-02-3}.

Cyclic refinement subsequently supplies a heat trace whose
limit is the theta series. The Mellin transform produces
$\pi^{-s/2}\Gamma(s/2)Z(s)$; scale inversion determines
the polar contribution and the functional equation of the
completed form. Propositions~\ref{rz-num-03-1}
and~\ref{rz-num-03-2} preserve the normalization and
convergence of this step. The finite-part and residual-readout
developments extend its evaluations with controlled
remainders. The Weil form jointly preserves the prime,
gamma and polar terms arising from this composition.

\subsection*{Double circle, reconstruction and transport}

The double circle distinguishes the spectral coordinate
from winding memory. The nine-phase decomposition produces
an exact balance with coefficient $8/81$; the central-block
modes determine the ratio $q=5/9$. The memory moments
$q^{|n|}$ have a positive triangular factorization in every
dimension. The loss register additionally preserves the
terminal term and cross products, so visible contraction
can be accompanied by isometric preservation of the
recorded history.

The analytic realization constructs a reversible operator
$C_a$ on exponentially decreasing test functions, with
$a>b>1/2$. Its integral action simultaneously preserves
correlations and polar cross products. The identity in
Subsection~\ref{tm-add:invariancia-weil} proves
\[
\mathscr W(C_af,C_ag)=\mathscr W(f,g).
\]
The margin $a>1/2$ belongs to the domain of this analytic
readout and ensures convergence of its polar contributions.
The regular memory circle and a representation receiving
an atomic limit retain their respective domains.

The gamma tail contains an explicit inverse readout:
it recovers the test function and reconstructs the negative
column from the positive one.
Equations~\eqref{gamma-add-eq-12}
and~\eqref{gamma-add-eq-14} describe its action, and
\eqref{gamma-add-eq-16} preserves the complete difference
of forms. This reconstruction provides the operator to be
studied; boundedness and contractivity are different
properties, whose estimates are developed next.

\subsection*{Coercivity and criterion for global extension}

Theorem~\ref{schur-add-schur-exacto} proves that, once
a positive bound for the detail block $D$ has been obtained,
positivity of the form on $V_R$ is equivalent to that of
$S=A-B^*D^{-1}B$. The correction retains the effect of
all details and their cross terms with the averages.
Theorem~\ref{schur-add-residual} bounds this matrix through
residuals and an operator error estimate. On the window
of length $1/9$, the estimates
$D>I$, $A>97/100$ and $\|B\|^2<1/5$ yield
Theorem~\ref{schur-add:coercividad-ventana}. The local
window satisfies $1/9<\log2$; its transported collections
may have tails and then receive the complete prime sum
through the proved invariance.

The global criterion also requires cross terms between
distinct collections and coverage of the entire test-function
domain. Its precise formulation is given in
Subsections~\ref{tm-add-m6-2}, \ref{tm-add-m7}
and~\ref{gamma-add-sub-375}: identification of the complete
arithmetic moments with a positive Gram form amounts to
joint preservation of energy and transport, with their
domains and limits. This condition remains a hypothesis
of the global results; proved coercivity belongs to the
explicit domains above. The article's contribution is to
connect arithmetic generation to reconstruction identities
and functional estimates precise enough to establish
that scope and its exact condition for extension.
\clearpage
